%; whizzy document
% latex beamer presentation.
% platex, latex-beamer でコンパイルすることを想定。 

%     Tokyo Debian Meeting resources
%     Copyright (C) 2006 Junichi Uekawa

%     This program is free software; you can redistribute it and/or modify
%     it under the terms of the GNU General Public License as published by
%     the Free Software Foundation; either version 2 of the License, or
%     (at your option) any later version.

%     This program is distributed in the hope that it will be useful,
%     but WITHOUT ANY WARRANTY; without even the implied warranty of
%     MERCHANTABILITY or FITNESS FOR A PARTICULAR PURPOSE.  See the
%     GNU General Public License for more details.

%     You should have received a copy of the GNU General Public License
%     along with this program; if not, write to the Free Software
%     Foundation, Inc., 51 Franklin St, Fifth Floor, Boston, MA  02110-1301 USA

% 実行順番
% sudo  ~/bin/usb-macbook-ir.c &
% real presentation (shell-command (concat "DISPLAY=:0.1 xpdf -fullscreen " (replace-regexp-in-string "tex$" "pdf"(buffer-file-name)) "&"))
% DISPLAY=:0.1 xpdf -fullscreen 

\documentclass[cjk,dvipdfmx]{beamer}
\usetheme{Warsaw}
%  preview (shell-command (concat "xpdf " (replace-regexp-in-string "tex$" "pdf"(buffer-file-name)) "&"))
%  presentation (shell-command (concat "xpdf -fullscreen " (replace-regexp-in-string "tex$" "pdf"(buffer-file-name)) "&"))

%http://www.naney.org/diki/dk/hyperref.html
%日本語EUC系環境の時
\AtBeginDvi{\special{pdf:tounicode EUC-UCS2}}
%シフトJIS系環境の時
%\AtBeginDvi{\special{pdf:tounicode 90ms-RKSJ-UCS2}}

\title{東京エリア Debian 勉強会}
\subtitle{資料}
\author{上川 純一 dancer@debian.org\\IRC nick: dancerj}
\date{2006年10月28日}
\logo{\includegraphics[width=8cm]{image200607/openlogo-light.eps}}

% 三択問題用
\newcounter{santakucounter}
\newcommand{\santaku}[5]{%
\addtocounter{santakucounter}{1}
\frame{\frametitle{問題\arabic{santakucounter}. #1}
%問題\arabic{santakucounter}. #1
\begin{minipage}[t]{0.7\hsize}
 \begin{itemize}
 \item  A #2\\
 \item  B #3\\
 \item  C #4\\
 \end{itemize}
\end{minipage}
}
\frame{\frametitle{問題\arabic{santakucounter}. #1}
%問題\arabic{santakucounter}. #1
\begin{minipage}[t]{0.7\hsize}
\begin{itemize}
\item  A #2\\
\item  B #3\\
\item  C #4\\
\end{itemize}
\end{minipage}
\begin{minipage}[t]{0.2\hsize}
答えは:


\vspace{1cm}

{\huge \hspace{1cm}#5}
\end{minipage}}
}


\begin{document}
\frame{\titlepage{}}

\section{Intro}

\begin{frame}
 \frametitle{本日のagenda}
\begin{minipage}[t]{0.4\hsize}
  \begin{itemize}
  \item 注意事項
	\begin{itemize}
	 \item 飲食禁止
	 \item トイレは？
	\end{itemize}
 \end{itemize}
\end{minipage} 
\begin{minipage}[t]{0.4\hsize}
 \begin{itemize}
  \item Debian 勉強会紹介
  \item 翻訳の仕組紹介
  \item IceWeaselとは何か？
 \end{itemize}
\end{minipage}
\end{frame}

\section*{講師紹介}
\begin{frame}{講師紹介：上川 純一}
 \begin{itemize}
  \item Debian Developer
  \item Debian JP 会長
  \item pbuilder, apt-listbugs, cowdancer などのパッケージメンテナンス
  \item Debian 勉強会首謀者
 \end{itemize}
\end{frame}

\section{Debian 勉強会とは}
\begin{frame}
 
「Debian力がつく会合です」

\hfill{}矢吹幸治
\end{frame}

\subsection{概要紹介}
\begin{frame}{Debian　勉強会概要紹介}
\begin{itemize}
 \item 経緯
       \begin{itemize}
	\item 2005年1月開始
	\item 月に一回勉強会
	\item 東京杉並区荻窪開催
	\item 15-20名程度の規模
       \end{itemize} 
 \item 目的
       \begin{itemize}
	\item IRC/ML だけではない会議の場所の提供
	\item ドキュメントの整備
	\item Debian Official Developerの輩出
       \end{itemize}
\end{itemize}
\end{frame}

\begin{frame}{内容例}
\begin{itemize}[<+->]
 \item Debianの日本語翻訳作業について
 \item MacBook に Debian インストール
 \item Debian Policy 勉強会
 \item LaTeX の利用方法
 \item Debian での Flash の作り方
 \item Debian のネットワーク設定についての議論
 \item 課題：　\\マニアックすぎるので、もうちょっと大衆受けする内容がほ
       しい？
\end{itemize}
\end{frame}

\section*{講師紹介}
\begin{frame}{講師紹介：小林儀匡}
 \begin{itemize}
  \item Debian GNU/Linux ユーザ
  \item skk関連のパッケージ管理
  \item Debian Weely News 翻訳者
  \item Debian JP 翻訳チームメンバー
 \end{itemize}
\end{frame}

\begin{frame}{講師紹介：やまねひでき}
 \begin{itemize}
  \item Debian GNU/Linux ユーザ
  \item ttf-vlgothic, eclipse-nls などパッケージメンテナンス
  \item Debian JP 翻訳チームメンバー
 \end{itemize}
\end{frame}


\section{今後の予定}
\begin{frame}{今後の予定}
\begin{itemize}
 \item 2006/11/19 関西開催
 \item 2006/12/16 一年のしめくくりの会議、杉並区荻窪
\end{itemize}

 参加おまちしております
\end{frame}



\end{document}
